\scanpage{24}
Ceci posé, regardons, pour un $G$ quasi-cohérent variable sur $S_0$, les $q_*\underline{\operatorname{Hom}}(\,,\unsure{\mathcal O_{\Gamma_0}|\underline{\mathcal O}_{\Gamma_0}}\otimes_{S_0}G)$, $q'_*\underline{\operatorname{Hom}}(\,,\mathcal O_{\Gamma'_0}\otimes_{S_0}G)$ relatives aux deux lignes, on trouve \inserted{(via les Modules représentatifs)} un diagramme commutatif de \deleted{suites exactes}

\diagram{\begin{tikzcd}[ampersand replacement=\&,column sep=normal]
\& \& L\arrow[r,"i"] \& M\arrow[r,"d"] \& N\arrow[r] \& 0\\
(D) \& 0\arrow[r] \& L'\arrow[u,"\alpha"']\arrow[r,"i'"] \& M'\arrow[u,"\beta"']\arrow[r,"d'"] \& N'\arrow[u,"\gamma"']\arrow[r,dash] \& 0\,, 
\end{tikzcd}}

de plus on a un hom can. $\sigma$

\diagram{\begin{tikzcd}[ampersand replacement=\&,row sep=small,column sep=small,execute at end picture={\draw ([xshift=-4pt,yshift=5pt]js.north west) -- ([xshift=-8pt,yshift=4pt]ls.north west) -- ([xshift=5pt,yshift=1pt]ls.north east); \draw ([xshift=-4pt,yshift=0pt]js.south west) -- ([yshift=-6pt]js.south) -- ([xshift=5pt,yshift=-1pt]ls.south east);}]
\& \& |[alias=ls]|L\arrow[dl,"\sigma"']\\
\& |[alias=js]|J\arrow[dl,dashed,"\lambda"'] \& L'\arrow[u,"\alpha"']\\
G
\end{tikzcd}}

\diagram{\begin{tikzcd}[ampersand replacement=\&]
L'\&M'\arrow[l,"\pi'"']
\end{tikzcd}}

correspondant à l'identité $\mathcal O_{\Gamma_0}\otimes_{S_0}J\longrightarrow\mathcal O_{\Gamma_0}\otimes_{S_0}\underset{\substack{\Vert\\G}}{J}$, et un \uline{splitting} \inserted{$M'\xrightarrow{\unsure{\pi|\pi'}}L'$} de la deuxième ligne.

Ceci dit, pour un $\lambda:J\to G$ donné, les solutions du pb relativement à $G$ sont les hom $M\xrightarrow{\mu}G$ qui relèvent \inserted{i.e. $\boxed{\mu i=\lambda\sigma}$} le composé $L\xrightarrow{\lambda\sigma}G$, et tels que $\mu\beta=(\underbrace{\lambda\sigma}\alpha)\pi'$. \deleted{\unsure{Il}} Pour qu'il en existe un, il f. et s. que

\margin{Condition de non ramification en $u_0$}

$1^\circ)$ $(\lambda\sigma)(\operatorname{Ker}i)=0$ i.e. $\lambda(R_1)=0$, où $R_1=\sigma(\operatorname{Ker}i)\subset J$

\margin{Alors on peut remplacer $L$ par $L/\operatorname{Ker}i$}

[\deleted{i.e. $\lambda$ se factorise en $\bar\lambda:\bar J=J/K\longrightarrow G$}].

$2^\circ)$ [Tenant compte que \uline{$\gamma$ \deleted{\unsure{est}}surjectif} ] $\lambda\sigma\alpha\pi'$ \unsure{vaut|nulle} sur \deleted{\unsure{$R'$}}$\operatorname{Ker}\beta$, i.e. $\lambda(R_2)=0$, où $R_2=\unsure{\lambda}\sigma\alpha\pi'(\operatorname{Ker}\beta)$

% archive sheet number: 23
